%==============================================================================
% Dàn ý chi tiết - Chapter 2: Cơ sở lý thuyết
%==============================================================================

\chapter{Cơ sở lý thuyết}

%-----------------------------------------------------------------------------
\section{Mô hình ngôn ngữ lớn (Large Language Models)}
%-----------------------------------------------------------------------------

\subsection{Kiến trúc Transformer}
\begin{itemize}
    \item Giới thiệu kiến trúc Transformer (Vaswani et al., 2017)
    \item Cơ chế Self-Attention:
    \begin{itemize}
        \item Query, Key, Value
        \item Multi-Head Attention
        \item Tính toán attention score
    \end{itemize}
    \item Kiến trúc Encoder-Decoder
    \begin{itemize}
        \item Position Encoding
        \item Feed Forward Networks
        \item Residual Connection và Layer Normalization
    \end{itemize}
    \item Ưu điểm của Transformer so với RNN/LSTM
    \item Các biến thể: Decoder-only (GPT), Encoder-only (BERT), Encoder-Decoder (T5, BART)
\end{itemize}

\subsection{Các mô hình ngôn ngữ lớn sử dụng trong nghiên cứu}
\begin{itemize}
    \item GPT Series (GPT-3.5, GPT-4):
    \begin{itemize}
        \item Đặc điểm kiến trúc
        \item Instruction Tuning và RLHF
        \end{itemize}
    \item Claude (Anthropic):
    \begin{itemize}
        \item Constitutional AI
        \item Haiku/Claude-3 Sonnet/Opus
    \end{itemize}
    \item LLaMA và các phiên bản open-source:
    \begin{itemize}
        \item LLaMA 2, LLaMA 3
        \item Mistral, Mixtral
    \end{itemize}
    \item So sánh các mô hình về khả năng sinh câu hỏi
    \item Các tham số quan trọng: temperature, top-p, max tokens
\end{itemize}

%-----------------------------------------------------------------------------
\section{Sinh câu hỏi trắc nghiệm tự động}
%-----------------------------------------------------------------------------

\subsection{Các phương pháp truyền thống}
\begin{itemize}
    \item Template-based Question Generation:
    \begin{itemize}
        \item Sử dụng câu mẫu có sẵn
        \item Thay thế từ khóa bằng placeholder
    \end{itemize}
    \item Rule-based approaches:
    \begin{itemize}
        \item Sử dụng ngữ pháp và cú pháp
        \item Named Entity Recognition để xác định câu trả lời
    \end{itemize}
    \item Statistical Machine Translation (SMT) based methods
    \item Hạn chế của phương pháp truyền thống:
    \begin{itemize}
        \item Thiếu tính đa dạng
        \item Chất lượng ngữ pháp kém
        \item Khó mở rộng
    \end{itemize}
\end{itemize}

\subsection{Sinh câu hỏi dựa trên mô hình ngôn ngữ lớn}
\begin{itemize}
    \item Prompt Engineering:
    \begin{itemize}
        \item Zero-shot prompting
        \item Few-shot prompting
        \item Chain-of-Thought (CoT) prompting
    \end{itemize}
    \item Fine-tuning LLMs cho question generation:
    \begin{itemize}
        \item Dataset preparation
        \item Loss function và training strategy
    \end{itemize}
    \item Các nghiên cứu liên quan:
    \begin{itemize}
        \item Self-QA (King et al.)
        \item SQuAD to MCQ conversion
        \item Các phương pháp sử dụng T5, BART
    \end{itemize}
    \item Thách thức:
    \begin{itemize}
        \item Sinh câu hỏi với độ khó đa dạng
        \item Tạo distractors chất lượng
        \item Đảm bảo answerability
    \end{itemize}
\end{itemize}

%-----------------------------------------------------------------------------
\section{Phân loại Bloom (Bloom's Taxonomy)}
%-----------------------------------------------------------------------------

\subsection{Cấp độ 1 -- Truy xuất (Retrieval)}
\begin{itemize}
    \item Định nghĩa và mục tiêu học tập
    \begin{itemize}
        \item Nhận biết, nhớ lại thông tin
        \item Xác định sự kiện, khái niệm cụ thể
    \end{itemize}
    \item Đặc điểm câu hỏi cấp độ 1:
    \begin{itemize}
        \item Câu hỏi "Ai?", "Khi nào?", "Ở đâu?"
        \item Yêu cầu nhắc lại thông tin có trong văn bản
    \end{itemize}
    \item Ví dụ câu hỏi cấp độ 1
    \item Chiến lược sinh câu hỏi Retrieval
\end{itemize}

\subsection{Cấp độ 2 -- Suy luận và tổng hợp (Inference \& Synthesis)}
\begin{itemize}
    \item Định nghĩa và mục tiêu học tập:
    \begin{itemize}
        \item Hiểu và diễn giải thông tin
        \item Tổng hợp kiến thức từ nhiều nguồn
        \item Đưa ra kết luận từ dữ liệu
    \end{itemize}
    \item Đặc điểm câu hỏi cấp độ 2:
    \begin{itemize}
        \item Câu hỏi "Tại sao?", "Làm thế nào?"
        \item Yêu cầu suy luận, so sánh, phân tích
    \end{itemize}
    \item Ví dụ câu hỏi cấp độ 2
    \item Chiến lược sinh câu hỏi Inference/Synthesis
\end{itemize}

\subsection{Cấp độ 3 -- Lập luận phản biện (Critical Reasoning)}
\begin{itemize}
    \item Định nghĩa và mục tiêu học tập:
    \begin{itemize}
        \item Đánh giá, phân tích phê phán
        \item Đưa ra quan điểm cá nhân
        \item Giải quyết vấn đề phức tạp
    \end{itemize}
    \item Đặc điểm câu hỏi cấp độ 3:
    \begin{itemize}
        \item Câu hỏi yêu cầu đánh giá, so sánh, phê phán
        \item Cần kiến thức ngoài văn bản
        \item Câu hỏi mở nhưng có thể trả lời bằng evidence
    \end{itemize}
    \item Ví dụ câu hỏi cấp độ 3
    \item Chiến lược sinh câu hỏi Critical Reasoning
    \item So sánh 3 cấp độ Bloom trong ngữ cảnh sinh câu hỏi
\end{itemize}

%-----------------------------------------------------------------------------
\section{Hệ thống đa tác tử (Multi-Agent Systems)}
%-----------------------------------------------------------------------------

\subsection{Khái niệm hệ thống đa tác tử}
\begin{itemize}
    \item Định nghĩa Agent trong AI
    \begin{itemize}
        \item Autonomous agents
        \item Reactive vs Proactive agents
    \end{itemize}
    \item Multi-Agent Systems (MAS):
    \begin{itemize}
        \item Định nghĩa và đặc điểm
        \item Ưu điểm so với Single-Agent
        \item Các loại tương tác: Cooperation, Competition, Negotiation
    \end{itemize}
    \item Ứng dụng của MAS trong NLP
    \begin{itemize}
        \item Debate systems
        \item Collaborative problem solving
        \item Code generation (Devin, SWE-Agent)
    \end{itemize}
    \item Lợi ích của Multi-Agent trong sinh câu hỏi:
    \begin{itemize}
        \item Phân tách nhiệm vụ rõ ràng
        \item Chuyên môn hóa từng agent
        \item Dễ dàng prompt engineering
    \end{itemize}
\end{itemize}

\subsection{Framework LangGraph}
\begin{itemize}
    \item Giới thiệu LangGraph:
    \begin{itemize}
        \item Khái niệm graph-based workflow
        \item So sánh với LangChain
    \end{itemize}
    \item Các thành phần cơ bản:
    \begin{itemize}
        \item Nodes (tác tử)
        \item Edges (luồng điều khiển)
        \item State management
    \end{itemize}
    \item Cycle và Conditional edges
    \begin{itemize}
        \item Xử lý vòng lặp
        \item Điều kiện rẽ nhánh
    \end{itemize}
    \item Persistence:
    \begin{itemize}
        \item Checkpointing
        \item Memory management
    \end{itemize}
    \item Ví dụ kiến trúc Multi-Agent với LangGraph
    \begin{itemize}
        \item State schema
        \item Agent definitions
        \item Workflow composition
    \end{itemize}
\end{itemize}

%-----------------------------------------------------------------------------
\section{Nhận dạng ký tự quang học (OCR)}
%-----------------------------------------------------------------------------

\subsection{OCR truyền thống}
\begin{itemize}
    \item Nguyên lý hoạt động của OCR truyền thống:
    \begin{itemize}
        \item Image preprocessing
        \item Text detection
        \item Character recognition
    \end{itemize}
    \item Các công cụ OCR phổ biến:
    \begin{itemize}
        \item Tesseract
        \item EasyOCR
        \item PaddleOCR
    \end{itemize}
    \begin{itemize}
        \item Ưu điểm: Nhanh, nhẹ, open-source
        \item Nhược điểm: Kém với font phức tạp, nhiễu
    \end{itemize}
    \item Các metrics đánh giá OCR:
    \begin{itemize}
        \item Character Error Rate (CER)
        \item Word Error Rate (WER)
    \end{itemize}
\end{itemize}

\subsection{OCR dựa trên mô hình thị giác--ngôn ngữ (Vision-Language Models)}
\begin{itemize}
    \item Giới thiệu VLM:
    \begin{itemize}
        \item CLIP
        \item GPT-4V
        \item Claude Vision
    \end{itemize}
    \item Phương pháp OCR với VLM:
    \begin{itemize}
        \item Prompt-based OCR
        \item Zero-shot text recognition
    \end{itemize}
    \item So sánh VLM OCR vs Traditional OCR:
    \begin{itemize}
        \item Độ chính xác trên ảnh phức tạp
        \item Khả năng xử lý đa ngôn ngữ
    \end{itemize}
    \item Hybrid approaches:
    \begin{itemize}
        \item Kết hợp Traditional OCR + VLM refinement
    \end{itemize}
    \item Mục đích OCR trong nghiên cứu:
    \begin{itemize}
        \item Tạo benchmark cho evaluation
        \item Mô phỏng scanning thực tế
    \end{itemize}
\end{itemize}

%-----------------------------------------------------------------------------
\section{Phương pháp đánh giá chất lượng câu hỏi}
%-----------------------------------------------------------------------------

\subsection{Đánh giá bởi con người}
\begin{itemize}
    \item Các tiêu chí đánh giá:
    \begin{itemize}
        \item Grammar and fluency
        \item Answerability
        \item Difficulty appropriateness
        \item Distractor plausibility
    \end{itemize}
    \begin{itemize}
        \item Relevance to passage
    \end{itemize}
    \item Thang điểm đánh giá (Likert scale)
    \begin{itemize}
        \item 1-5 scale
        \item Definitions for each score level
    \end{itemize}
    \begin{itemize}
        \item Inter-annotator agreement (Cohen's Kappa)
    \end{itemize}
    \begin{itemize}
        \item Hạn chế: Tốn kém, không thể mở rộng
    \end{itemize}
\end{itemize}

\subsection{LLM-as-a-Judge}
\begin{itemize}
    \item Khái niệm và nguyên lý:
    \begin{itemize}
        \item Sử dụng LLM để đánh giá đầu ra của LLM
    \end{itemize}
    \begin{itemize}
        \item Prompt engineering cho evaluation prompts
    \end{itemize}
    \begin{itemize}
        \item Scoring rubrics
    \end{itemize}
    \begin{itemize}
        \item Chain-of-thought evaluation
    \end{itemize}
    \item Các nghiên cứu liên quan:
    \begin{itemize}
        \item ChatArena
        \item AlpacaEval
        \item MT-Bench
    \end{itemize}
    \begin{itemize}
        \item Ưu điểm: Scalable, consistent, cost-effective
    \end{itemize}
    \begin{itemize}
        \item Nhược điểm: Potential bias, self-preference
    \end{itemize}
    \begin{itemize}
        \item Best practices: Pairwise comparison, rubric-based scoring
    \end{itemize}
\end{itemize}
